\section{Design Lessons by Pillar}
\label{sec:pillar-lessons}

Section~\ref{sec:objectives-retrospective} summarises the evaluation outcomes. This section 
serves as a record for the design lessons from building and testing MERIT. It covers how the implementation architecture should be 
understood, and also the where evaluation exposed the gaps that will be addressed later in Section~\ref{sec:future-work-roadmap}.
It is not a list of features, but rather a record of design decisions and the reasoning behind such choices.

\subsection{Pillar 1: Treat ingestion as lossy}
\label{sec:pillar1-reflection}

As we saw in Study~04, there are two main sources of loss when it comes to ingestion. Firstly, a parser
that fails to account for variances in CV formats will often fail to provide a fair and thorough assessment of a candidate
when they later get scored against a job description. This is evidently noticed in the difference in parsing
for names and experience between single-column and multi-column layouts. Secondly, a dictionary that is not comprehensive 
enough will also fail to extract all the skills that a candidate has to offer, once again, unfairly penalising candidates
that are not on the basis of merit. LinkedIn can help serve as a fallback for potentially missed CV experience while GitHub
can help serve as a fallback for potentially missed CV skills.

\subsection{Pillar 2: Missing data $\neq$ low skill}
\label{sec:pillar2-reflection}

The lesson here is not that multi-source fusion is missing, but rather that the current scoring policy
treats absent GitHub activity as weak/no evidence, rather than a potential signal of missing data. This is
evidenced in Study~01B, where Carl Corp is penalised for empty GitHub despite significant tenure.
Research into missing data policies is therefore a promising direction for future work.

\subsection{Pillar 3: Explain conflict, not only attribution}
\label{sec:pillar3-reflection}
Glass-box output is core to MERIT's design, from the implementation of Shapley values to the integration of Blind Mode.
Negative source contributions and integrity-driven reversals show the Glass-Box working as intended. However, the key
lesson we need to note is that we must understand the limits of what was shipped. Study~10 shows that the system can still
accompany a wrong rank unless we are able to surface dissonance appropriately.
